% Options for packages loaded elsewhere
\PassOptionsToPackage{unicode}{hyperref}
\PassOptionsToPackage{hyphens}{url}
\documentclass[
]{article}
\usepackage{xcolor}
\usepackage{amsmath,amssymb}
\setcounter{secnumdepth}{-\maxdimen} % remove section numbering
\usepackage{iftex}
\ifPDFTeX
  \usepackage[T1]{fontenc}
  \usepackage[utf8]{inputenc}
  \usepackage{textcomp} % provide euro and other symbols
\else % if luatex or xetex
  \usepackage{unicode-math} % this also loads fontspec
  \defaultfontfeatures{Scale=MatchLowercase}
  \defaultfontfeatures[\rmfamily]{Ligatures=TeX,Scale=1}
\fi
\usepackage{lmodern}
\ifPDFTeX\else
  % xetex/luatex font selection
\fi
% Use upquote if available, for straight quotes in verbatim environments
\IfFileExists{upquote.sty}{\usepackage{upquote}}{}
\IfFileExists{microtype.sty}{% use microtype if available
  \usepackage[]{microtype}
  \UseMicrotypeSet[protrusion]{basicmath} % disable protrusion for tt fonts
}{}
\makeatletter
\@ifundefined{KOMAClassName}{% if non-KOMA class
  \IfFileExists{parskip.sty}{%
    \usepackage{parskip}
  }{% else
    \setlength{\parindent}{0pt}
    \setlength{\parskip}{6pt plus 2pt minus 1pt}}
}{% if KOMA class
  \KOMAoptions{parskip=half}}
\makeatother
\usepackage{longtable,booktabs,array}
\usepackage{caption}
\captionsetup[table]{skip=6pt}
\newcounter{none} % for unnumbered tables
\usepackage{calc} % for calculating minipage widths
% Correct order of tables after \paragraph or \subparagraph
\usepackage{etoolbox}
\makeatletter
\patchcmd\longtable{\par}{\if@noskipsec\mbox{}\fi\par}{}{}
\makeatother
% Allow footnotes in longtable head/foot
\IfFileExists{footnotehyper.sty}{\usepackage{footnotehyper}}{\usepackage{footnote}}
\makesavenoteenv{longtable}
\setlength{\emergencystretch}{3em} % prevent overfull lines
\providecommand{\tightlist}{%
  \setlength{\itemsep}{0pt}\setlength{\parskip}{0pt}}
\usepackage{bookmark}
\IfFileExists{xurl.sty}{\usepackage{xurl}}{} % add URL line breaks if available
\urlstyle{same}
% fallback for those not using the hyperref driver hyperxmp:
\makeatletter
\@ifundefined{xmpquote}{\newcommand{\xmpquote}[1]{#1}}{}
\makeatother
\hypersetup{
  hidelinks,
  pdfcreator={LaTeX via pandoc}}

\author{}
\date{}

\begin{document}

\section{Informe de puerta E6 --- run\_v1.22\_r3 (Paper IV
v1.22-r3)}\label{informe-de-puerta-e6--run_v122_r3-paper-iv-v122-r3}

Veredicto actual: \textbf{PASS\_WITH\_FINDINGS (NEW-05)}. Evidencia:
\texttt{20\_EVIDENCE/E6/}. Auditor: Claude Opus 5.5, misma sesión que
las ejecuciones anteriores; no se ejecutó Lean.

\begin{center}\rule{0.5\linewidth}{0.5pt}\end{center}

\section{E6 --- NEW-04 and parity
(run\_v1.22\_r3)}\label{e6--new-04-and-parity-run_v122_r3}

\subsection{Method (independent of the editor\textquotesingle s
scripts)}\label{method-independent-of-the-editors-scripts}

\begin{itemize}
\tightlist
\item
  \textbf{Diff and protected content.}
  \texttt{10\_SCRIPTS/e6\_r3\_delta.py} produces
  \texttt{E6\_R3\_DELTA.json} and the auditor\textquotesingle s diffs
  \texttt{auditor\_diff\_\{en,es\}\_\{md,tex\}.diff}.

  \begin{itemize}
  \tightlist
  \item
    Paragraphs: 503 → 503 in each language; \textbf{5 changed paragraph
    groups per language}.
  \item
    Unchanged in EN and ES: 237 displays, 1989 inline formulas, code
    blocks, 146 tags, 159 table rows, headings, image references and the
    whole bibliography.
  \item
    Changed paragraphs contain no inline mathematics in either r2 or r3.
  \item
    Code spans: 331 → 333. No span is removed; the two added are
    \texttt{run\_v1.22\_r1} and \texttt{run\_v1.22\_r2} (audit-run
    names, not Lean identifiers).
  \item
    TeX: preamble identical; 10 changed lines per language; displays and
    equation environments identical.
  \item
    The body of the editor\textquotesingle s
    \texttt{CHANGES\_\{en,es\}.diff} equals the
    auditor\textquotesingle s diff byte for byte. Only the two header
    lines (file names) differ.
  \end{itemize}
\item
  \textbf{ES bold leads.} The single difference among bold paragraph
  leads is the A.2 lead, «Alcance de la demostración escrita» → «Alcance
  de la prueba escrita». It sits inside a changed paragraph and is not a
  section heading.
\item
  \textbf{PDF corresponds to TeX.}

  \begin{itemize}
  \tightlist
  \item
    EN r3: 72 pages, xdvipdfmx 2026-10-01 13:56 UTC. ES r3: 73 pages,
    13:54 UTC.
  \item
    Every token-diff region between the r2 and r3 PDF texts (EN 56, ES
    63) is explained by the changed MD paragraphs, with 0 unexplained
    regions.
  \end{itemize}
\item
  \textbf{Raster (72 dpi).} EN changes on pp. 1, 33, 39, 40, 64; ES on
  pp. 1--6, 34--36, 39, 40, 66. This equals
  \texttt{ARTIFACT\_CHECKS.json}.
\item
  \textbf{Visual inspection} (\texttt{PAGE\_INSPECTION\_LOG.csv},
  145/145).

  \begin{itemize}
  \tightlist
  \item
    Every page was seen on 4-up sheets at 70 dpi (\texttt{overview/}).
  \item
    The pages with changes or shifts were also inspected in
    r2\textbar r3 side-by-side comparisons at 105 dpi (\texttt{cmp/}):
    ES 36, 39, 40, 66 and shifts.
  \item
    For pages whose comparison images could not be confirmed, the rows
    were re-logged from displayed sheets (CORRECTIONS C4-01).
  \item
    The ES shift (pp. 1--6, 34--36) is absorbed by white space before
    Figure 1 and Figure 4, so pp. 7--33 and 37+ are raster-identical. No
    content is lost or duplicated, and page counts are unchanged.
  \end{itemize}
\end{itemize}

\subsection{NEW-04 --- the five documentary paragraphs per
language}\label{new-04--the-five-documentary-paragraphs-per-language}

All ten paragraphs (five EN, five ES) are quoted in full before/after in
\texttt{E6\_R3\_DELTA.json} (\texttt{changed}).

{\def\LTcaptype{none} % do not increment counter
\begin{longtable}[]{@{}llll@{}}
\toprule\noalign{}
Place & r2 (stale) & r3 & Check \\
\midrule\noalign{}
\endhead
\bottomrule\noalign{}
\endlastfoot
Front matter (Status/Estado) & «revisión que responde a la revalidación
de 1.21 \ldots{} INCONCLUSIVE \ldots{} requieren su propia revisión» &
r2 concluded PASS\_WITH\_OBSERVATIONS on 1 Oct 2026; accepted C.3 and
the seven A.2 rows; no open finding ≥ MINOR; Lean cut unchanged with
formal PASS inherited; r3 updates the status and its delta awaits
verification; no new build & \textbf{Correct.} It does not invent a
clean PASS and it states the observations. \\
§7 «Audit status / Estado de auditoría» & «v1.22 \ldots{} esperan
revalidación» & v1.2 and v1.21 INCONCLUSIVE; C.3 accepted in
run\_v1.22\_r1; r2 PASS\_WITH\_OBSERVATIONS, observations retained; no
recompilation; r3 claims no verdict for its own changes &
\textbf{Correct}, and the historical INCONCLUSIVE verdicts are kept. \\
A.1, note after Table 7 & «La comparación ampliada de C.3 en v1.22
espera una nueva revisión» & evidence revalidated without recompilation;
r2 records C.3 and distinguishes individual semantic rows from coverage
by the component ledger and compilation & \textbf{Correct.} This is
consistent with run\_v1.22\_r2 §4.1. \\
A.2 «Scope of the written proof / Alcance de la prueba escrita» &
«aceptó seis de las siete filas \ldots{} requiere revisión» & the seven
rows are ACCEPTABLE\_SUMMARY per r2; R-01 limit retained; build ≠
exposition ≠ peer review & \textbf{Correct.} The A.2 table header
«Derivation to review\ldots/Derivación por revisar\ldots» is unchanged
because table rows are protected. It reads as guidance for a referee and
is acceptable as is. \\
F.4, closing paragraph & «Una nueva auditoría todavía debe
comprobar\ldots» & the E.4/F.3/F.5 checks are consolidated in r2, with
its scope; compilation is not a substitute for review &
\textbf{Correct.} run\_v1.2\_r1 E2 §5.1 rederived E.4 block by block,
F.3a and F.5 cases. \\
\end{longtable}
}

\textbf{C.3 and the seven A.2 rows are no longer presented as pending in
either language. No clean PASS is invented. EN/ES content parity holds.}

\textbf{NEW-04: CLOSED in r3.} The paragraphs are dated records of r2;
the verdict on r3 is this report. That is not an inconsistency and does
not require editing the manuscript retroactively (mandate §2).

\subsection{NEW-05 --- ES terminology in the r3 rewrites
(MINOR)}\label{new-05--es-terminology-in-the-r3-rewrites-minor}

The new ES text of the A.2 and F.4 paragraphs introduces terms that
differ from the established terminology of the ES manuscript:

\begin{itemize}
\tightlist
\item
  «la \textbf{programación numérica} de C.3» (A.2, MD l.1355). Elsewhere
  the text says «calendario numérico» (l.1599), the C.3 heading is «Un
  \textbf{calendario} explícito y su propagación», and r2 itself said
  «el calendario numérico en C.3».
\item
  «la cota adaptada \textbf{con muestra fijada}» (F.4, l.2356). F.5 says
  «Ensamblaje de la cota adaptada para \textbf{muestras con anclajes}»
  (l.2403), and Table 10 says «Cota de muestra con anclajes» (l.2398);
  r2 said «muestras con anclajes».
\item
  (Checked and \emph{not} a finding: «conteos de normalización» replaces
  r2\textquotesingle s «cuentas de normalización», but «Conteos» is
  established in the A.2 table, as in «Conteos de descarte».)
\end{itemize}

These are not mathematical changes, and EN/ES parity of meaning holds.
However, they break the terminological consistency the series maintains
(compare X-04 and X-15). A reader of the A.2 and F.4 notes might not
identify the named objects.

\textbf{Action before publication:} restore «calendario numérico» (A.2)
and «muestras con anclajes» (F.4) in the ES text, then repeat the ES
identity and E6 check.

\subsection{E6 verdict (r3)}\label{e6-verdict-r3}

\textbf{PASS\_WITH\_FINDINGS.}

\begin{itemize}
\tightlist
\item
  NEW-04 is closed; the protected content is identical; each PDF
  corresponds to its TeX; all 145 pages were inspected.
\item
  NEW-05 (MINOR, ES terminology) remains open, with a small
  pre-publication correction.
\end{itemize}

\end{document}
